% DIAPOSITIVA 19: Pool de 9 variables y selección del modelo
\begin{frame}{Modelado: feature engineering y selección del modelo}
  \begin{columns}[T]
    \begin{column}{0.50\textwidth}
      \textbf{Pool de 9 variables} (tras filtro VIF sobre 22 candidatas)
      \begin{enumerate}\small
        \item Logit dominancia de \textit{likes}
        \item Logit dominancia de comentarios
        \item Logit dominancia de compartidos
        \item Logit \textit{likes} por publicación
        \item Logit comentarios por publicación
        \item Logit compartidos por publicación
        \item Log-población del municipio
        \item Número de candidatos en la contienda
        \item Moderador: (1) por (log-pob.\ centrada)
      \end{enumerate}
      \vspace{0.2em}
      {\small\textbf{ElasticNet} ($L_1 + L_2$): selección automática.\\
      $\hat\alpha=0{,}0475$, $\lambda=0{,}10$, StandardScaler.}
    \end{column}
    \begin{column}{0.46\textwidth}
      \textbf{Selección del modelo (AutoML)}
      \begin{itemize}\small
        \item Familias evaluadas: GLM fraccional, gradient boosting, redes neuronales,
              SVM, random forests
        \item Criterio único: MAE fuera de muestra (LOCO-CV)
        \item \kw{Ganador: GLM fraccional + ElasticNet}
      \end{itemize}
      \vspace{0.3em}
      \begin{exampleblock}{¿Por qué GLM sobre gradient boosting?}
        \scriptsize (1) Cuota $\in(0,1)$: GLM binomial+logit garantiza $\hat{y}\in(0,1)$
        (Papke \& Wooldridge, 1996). (2) Interpretabilidad: coeficientes en log-odds.
        MAE similar (9,56 vs.\  9,91 pp).
      \end{exampleblock}
    \end{column}
  \end{columns}
\end{frame}

% DIAPOSITIVA 20: LOCO-CV
\begin{frame}{Modelado: validación LOCO-CV \corr{[Corr.\ 2.3]}}
  \begin{columns}[T]
    \begin{column}{0.52\textwidth}
      \textbf{Dos fases del proceso:}
      \begin{itemize}\small
        \item \textbf{Fase 1 (evaluación):} ciclo de 31 pliegues. En cada uno se bloquea
              una ciudad completa, se hace grid search interno (120 combinaciones de
              $\hat\alpha$ y $\lambda$) sobre las 30 ciudades restantes, se elige la
              mejor combinación, se entrena el ElasticNet y se predice la ciudad
              bloqueada. Resultado: 31 errores cuyo promedio es el MAE.
        \item \textbf{Fase 2 (modelo final):} con la combinación ganadora
              ($\hat\alpha=0{,}0475$, $\lambda=0{,}10$) se entrena el ElasticNet sobre
              las 31 ciudades completas. Los betas resultantes son los definitivos.
      \end{itemize}
      \vspace{0.2em}
      \begin{alertblock}{\small Por qué no CV estándar}
        \small Mezclar candidatos entre ciudades genera fuga de información:
        las cuotas de una contienda suman 1, por lo que ver un candidato de
        Medellín en entrenamiento adelanta información sobre los demás.
      \end{alertblock}
    \end{column}
    \begin{column}{0.44\textwidth}
      \begin{exampleblock}{Resultado LOCO-CV}
        \large
        \[\text{MAE}_{\text{oos}} = \mathbf{9{,}56\text{ pp}}\]
        \[R^2_{\text{oos}} = \mathbf{0{,}518}\]
      \end{exampleblock}
      \vspace{0.3em}
      \begin{block}{Grid search anidado}
        \small El CV interno elige $\hat\alpha$ y $\lambda$ \textbf{sin ver}
        la ciudad de prueba. La Fase 1 mide el MAE; la Fase 2 entrega los
        betas para aplicar en 2026.
      \end{block}
    \end{column}
  \end{columns}
\end{frame}

% DIAPOSITIVA 21: Especificación ganadora e interpretación de los betas
\begin{frame}{Modelado: especificación ganadora e interpretación}
  \textbf{ElasticNet retuvo \kw{2 variables} del pool de 9, las 7 restantes penalizadas a cero:}
  \[
    \hat{y}_c = \underbrace{0{,}258}_{\beta_0}
    + \underbrace{0{,}111}_{\beta_1} \cdot f_{\ell}
    + \underbrace{0{,}026}_{\beta_2} \cdot f_{\ell} \cdot (\ln\!P - \overline{\ln P})
  \]
  {\small $f_{\ell}$ = logit(dominancia de \textit{likes}); $P$ = población municipal}

  \vspace{0.3em}
  \begin{columns}[T]
    \begin{column}{0.50\textwidth}
      \textbf{Interpretación de los coeficientes:}
      \begin{itemize}\small
        \item $\beta_0 = 0{,}258$: cuota del candidato promedio en ciudad promedio
              ($\approx 1/3{,}84$ candidatos)
        \item $\beta_1 = 0{,}111$: por cada unidad adicional en logit(dominancia),
              la cuota aumenta aprox.\ 10 pp
        \item $\beta_2 = 0{,}026$: en ciudades grandes la señal digital
              pesa más electoralmente
      \end{itemize}
    \end{column}
    \begin{column}{0.46\textwidth}
      \begin{exampleblock}{\textbf{P3}: Respondida}
        \small La dominancia relativa de \textit{likes} en logit es el predictor
        más estable y robusto. Población modera el efecto (no predice directamente).
      \end{exampleblock}
      \vspace{0.3em}
      {\small\textbf{Criterio único de selección:}
      MAE fuera de muestra (LOCO-CV).\\
      Ninguna decisión tomada por ajuste en entrenamiento.}
    \end{column}
  \end{columns}
\end{frame}

% DIAPOSITIVA 20 (renumerada): Los dos modelos
\addtocounter{framenumber}{1}

% DIAPOSITIVA nueva: Estabilidad de coeficientes (bootstrap LOCO-CV)
\addtocounter{framenumber}{1}

% DIAPOSITIVA nueva: Diagnóstico LOCO-CV (predicción vs. real + residuos)
\addtocounter{framenumber}{1}
